\usepackage{pgf}
\usepackage{layout}
\usepackage{fancyhdr}
\usepackage{imakeidx}
\usepackage{etoolbox}

\usepackage{lipsum}

\usepackage{titlesec}
\usepackage{afterpage}
\usepackage{xfrac}
\usepackage{xcolor,sectsty}
\usepackage{gensymb}
\usepackage[most]{tcolorbox}

\usepackage{textgreek}
\usepackage{svg}

\usepackage{tikz}
\usepackage{multicol}
\usepackage{wrapfig}
\usepackage{caption}

\usepackage[T1]{fontenc}
\usepackage{kmath}
\usepackage{kerkis}

\usepackage{hyperref}

\usetikzlibrary{shapes.geometric}
\usetikzlibrary{decorations.pathmorphing}
\usetikzlibrary{calc}

\def\IsInteger#1{%
  TT\fi
  \begingroup \lccode`\-=`\0 \lccode`+=`\0
    \lccode`\1=`\0 \lccode`\2=`\0 \lccode`\3=`\0
    \lccode`\4=`\0 \lccode`\5=`\0 \lccode`\6=`\0
    \lccode`\7=`\0 \lccode`\8=`\0 \lccode`\9=`\0
  \lowercase{\endgroup
    \expandafter\ifx\expandafter\delimiter
    \romannumeral0\string#1}\delimiter
}

\makeatletter
\def\closeopenmulticols{%
% test if current env is "multicols" if so close it
   \def\@tempa{multicols}%
   \ifx\@tempa\@currenvir
      \end{multicols}%
  \fi 
}
\makeatother

\newlength{\currentfontsize}
\makeatletter
\newcommand{\refreshcurrentfontsize}{
	\setlength{\currentfontsize}{\f@size pt}
}
\makeatother

\rmfamily

%Steven Segletes's Faux Small Caps command for situations where bolded small caps are needed.
\newcommand\fauxsc[1]{\fauxschelper#1 \relax\relax}
\def\fauxschelper#1 #2\relax{%
	\fauxschelphelp#1\relax\relax%
	\if\relax#2\relax\else\ \fauxschelper#2\relax\fi%
}
\def\Hscale{.85}\def\Vscale{.72}\def\Cscale{1.10}
\def\fauxschelphelp#1#2\relax{%
	\ifnum`#1>``\ifnum`#1<`\{\scalebox{\Hscale}[\Vscale]{\uppercase{#1}}\else%
	\scalebox{\Cscale}[1]{#1}\fi\else\scalebox{\Cscale}[1]{#1}\fi%
	\ifx\relax#2\relax\else\fauxschelphelp#2\relax\fi}

%% COLOR SET UP STUFF
%
\newcounter{partpagedetector}
\AddToHook{cmd/part/before}{\setcounter{partpagedetector}{1}}
\AddToHook{cmd/chapter/before}{\setcounter{partpagedetector}{0}}

\definecolor{grptitle}{RGB}{200,212,220}
\definecolor{grpbgline}{RGB}{180,180,180}

\definecolor{chapterbgcolor1}{HTML}{C1FF85}
\definecolor{chapterbgcolor2}{HTML}{85FFFF}
\definecolor{chapterbgcolor3}{HTML}{C285FF}
\definecolor{chapterbgcolor4}{HTML}{FF8585}
\definecolor{chapterbgcolor5}{HTML}{FFED85}
\definecolor{chapterbgcolor6}{HTML}{85FFB0}
\definecolor{chapterbgcolor7}{HTML}{8597FF}
\definecolor{chapterbgcolor8}{HTML}{FF85D4}
\definecolor{chapterbgcolor9}{HTML}{85FFDC}
\definecolor{chapterbgcolor10}{HTML}{9F85FF}
\definecolor{chapterbgcolor11}{HTML}{FF85A8}
\definecolor{chapterbgcolor12}{HTML}{E5FF85}
\definecolor{chapterbgcolordefault}{HTML}{AAAAAA}

\newcommand{\ChapterBackgroundTint}{chapterbgcolordefault}

\def\SetChapterTint{%
	\ifnum\arabic{chapter}>12%
		\renewcommand{\ChapterBackgroundTint}{chapterbgcolordefault}
	\else%
		\ifnum\arabic{chapter}<1%
			\renewcommand{\ChapterBackgroundTint}{chapterbgcolordefault}
		\else%
			\edef\ChapterBackgroundTint{chapterbgcolor\arabic{chapter}}
		\fi%
	\fi%
	\ifnum\value{partpagedetector}=1
		\renewcommand{\ChapterBackgroundTint}{chapterbgcolordefault}
	\fi
}

%Command shortcut for GammaRP with greek gamma
\newcommand{\GammaRP}{\textGamma RP}

%Command shortcut for tikz dodecahedron
\pgfkeys{
	/GammaDodec Options/.cd,
	color/.initial={grpbgline},
	size/.initial={\currentfontsize},
}
\newlength{\gammadodecsize}
\newcommand{\GammaDodec}[2][]{%
\pgfkeys{/GammaDodec Options/.cd,#1}%
\refreshcurrentfontsize%
\setlength{\gammadodecsize}{\pgfkeysvalueof{/GammaDodec Options/size}}
\begin{tikzpicture}[remember picture, baseline=-0.3\gammadodecsize, line cap=round,line width=0.08\gammadodecsize,line join=round,draw=\pgfkeysvalueof{/GammaDodec Options/color}, inner sep=-4\gammadodecsize, outer sep=-4\gammadodecsize]
	\foreach \angle in
		{18, 90, 162, 234, 306}
	{
		\draw[line width=0.08\gammadodecsize] (\angle:0.43\gammadodecsize) -- (\angle:0.59\gammadodecsize);
		\draw[line width=0.08\gammadodecsize] (\angle:0.42\gammadodecsize) -- (\inteval{\angle+72}:0.42\gammadodecsize);
		\draw[line width=0.08\gammadodecsize] (\angle:0.6\gammadodecsize) -- (\inteval{\angle+36}:0.58\gammadodecsize);
		\draw[line width=0.08\gammadodecsize] (\inteval{\angle+36}:0.58\gammadodecsize) -- (\inteval{\angle+72}:0.6\gammadodecsize);
	}
	\node[rectangle] at (0,0) {#2};
\end{tikzpicture}%
{}}

%Dice commands
\newcommand{\GRPitalNonInt}[1]{%
 \if\IsInteger{#1}%
 	#1%
 \else%
 	\textit{#1}%
 \fi%
}

\newcommand{\GRProll}[2]{\mbox{%
	\if\IsInteger{#2}%
		\ifnum#2<0%
			D\GRPitalNonInt{#1}d\inteval{-1*(#2)}
		\else%
			D\GRPitalNonInt{#1}a#2
		\fi%
	\else%
		D\textit{#1}a\textit{#2}
	\fi%
}}

\newcommand{\GRProlla}[2]{\mbox{%
	D\GRPitalNonInt{#1}a\GRPitalNonInt{#2}
}}

\newcommand{\GRProlld}[2]{\mbox{%
	D\GRPitalNonInt{#1}d\GRPitalNonInt{#2}
}}


%Formatting a stat
\newcommand{\stat}[1]{ \{ \textsc{#1}  \} }

%Page background commands
\newlength{\currenttop}%to calc page margins for reduced complexity
\newlength{\currentleft}%
\newlength{\currentright}%
\newlength{\currentbottom}%
\newlength{\currenthtop}%header top
\newlength{\currenthbase}%header base
\newlength{\extragradlength}% how much the gradient should run into the text area
\newlength{\currentengr}%angled line border extent

\newcommand{\GRPpageBG}{%
	\setlength{\currenttop}{\dimeval{1in + \voffset +\headsep + \topmargin + \headheight} }
	\setlength{\currentbottom}{\dimeval{\paperheight - \currenttop - \textheight} }
	\ifodd\value{page}
		\setlength{\currentleft}{\dimeval{1in + \hoffset + \oddsidemargin} }
	\else
		\setlength{\currentleft}{\dimeval{1in + \hoffset + \evensidemargin} }
	\fi
	\setlength{\currentright}{\dimeval{\paperwidth - \currentleft - \textwidth} }
	\setlength{\extragradlength}{1.0cm}
	\setlength{\currentengr}{\headheight+3mm}
	\begin{tikzpicture}[remember picture,overlay]
	\node[anchor = north west, shift={(-3.7mm,4mm)}] at (current page.north west) {%
	\begin{tikzpicture}
		\shade[shading = axis, left color = \ChapterBackgroundTint, right color = \ChapterBackgroundTint!0] (-2mm,0) rectangle (\currentleft+\extragradlength, -\paperheight);
		\shade[shading = axis, right color = \ChapterBackgroundTint, left color = \ChapterBackgroundTint!0] (\paperwidth-\currentright-\extragradlength, 0) rectangle (\paperwidth+2mm,-\paperheight);
		\shade[shading = axis, top color = \ChapterBackgroundTint, bottom color = \ChapterBackgroundTint!0] (-2mm,2mm) -- (\currentleft+\extragradlength, -\currenttop-\extragradlength) -- (\paperwidth-\currentright-\extragradlength, -\currenttop-\extragradlength) -- (\paperwidth+2mm,2mm) -- cycle;
		\shade[shading = axis, bottom color = \ChapterBackgroundTint, top color = \ChapterBackgroundTint!0] (-2mm,-\paperheight-2mm) -- (\currentleft+\extragradlength, \currenttop+\extragradlength-\paperheight) -- (\paperwidth-\currentright-\extragradlength, \currenttop+\extragradlength-\paperheight) -- (\paperwidth+2mm,-\paperheight-2mm) -- cycle;	

		\foreach \vertpos in 
			{0.1,0.15,0.2,0.25,0.3,0.35,0.4,0.45,0.5,0.55,0.6,0.65,0.7,0.75,0.8,0.85,0.9}
		{%
			\shade[shading=axis, left color=\ChapterBackgroundTint!50!white, middle color=white, right color=white, shift={(0,-\vertpos\paperheight)}, decorate, decoration={zigzag,segment length=0.2cm}]  (\currentleft+\extragradlength+1mm,-0.4cm) -- (0,0) -- (\currentleft+\extragradlength+1mm,0.4cm);
			\shade[shading=axis, right color=\ChapterBackgroundTint!50!white, middle color=white, left color=white, shift={(\paperwidth,-\vertpos\paperheight)}, decorate, decoration={zigzag,segment length=0.2cm}]  (-\currentright-\extragradlength-1mm,0.4cm) -- (0,0) -- (-\currentright-\extragradlength-1mm,-0.4cm);
		}
		\foreach \horzpos in
			{0.15,0.25,0.35,0.45,0.55,0.65,0.75,0.85}
		{%
			\shade[shading=axis, top color=\ChapterBackgroundTint!50!white, middle color=white, bottom color=white, shift={(\horzpos\paperwidth,0)}, decorate, decoration={zigzag,segment length=0.2cm}] (-0.4cm,-\currenttop-\extragradlength-1mm) -- (0,0) -- (0.4cm,-\currenttop-\extragradlength-1mm);
			\shade[shading=axis, bottom color=\ChapterBackgroundTint!50!white, middle color=\ChapterBackgroundTint!0, top color=\ChapterBackgroundTint!0, shift={(\horzpos\paperwidth,-\paperheight)}, decorate, decoration={zigzag,segment length=0.2cm}] (0.4cm,\currentbottom+\extragradlength+1mm) -- (0,0) -- (-0.4, \currentbottom+\extragradlength+1mm);
		}

		\draw[draw=\ChapterBackgroundTint!90!black, line width=1mm] (\currentengr-1mm,1mm) -- (\currentengr*2,-\currentengr) -- (\paperwidth-\currentengr*2,-\currentengr) -- (\paperwidth-\currentengr+1mm,1mm);
		\draw[draw=\ChapterBackgroundTint!90!black, line width=1mm] (\currentengr-1mm,-\paperheight-1mm) -- (\currentengr*2,\currentengr-\paperheight) -- (\paperwidth-\currentengr*2,\currentengr-\paperheight) -- (\paperwidth-\currentengr+1mm,-\paperheight-1mm);
		\draw[draw=\ChapterBackgroundTint!90!black, line width=1mm] (-1mm,1mm-\currentengr) -- (\currentengr,-\currentengr*2) -- (\currentengr,-\paperheight+\currentengr*2) -- (-1mm,-\paperheight+\currentengr-1mm);
		\draw[draw=\ChapterBackgroundTint!90!black, line width=1mm] (\paperwidth+1mm,1mm-\currentengr) -- (\paperwidth-\currentengr,-\currentengr*2) -- (\paperwidth-\currentengr,-\paperheight+\currentengr*2) -- (\paperwidth+1mm,-\paperheight+\currentengr-1mm);

		\foreach \currshiftx / \currshifty / \currrot in { 0/0/0, {\paperwidth}/0/-90, 0/{-\paperheight}/90, {\paperwidth}/{-\paperheight}/180}
		{%
			\draw[draw=\ChapterBackgroundTint!90!black, line width=1mm, shift={(\currshiftx,\currshifty)}, rotate=\currrot] (-0.2,0cm) -- (0.6cm,-0.8cm) -- (0.7cm,-1.0cm) -- (0.8cm,-1cm) -- (1cm,-1cm) -- (1cm,-0.8cm) -- (1.0cm,-0.7cm) -- (0.8cm, -0.6cm) -- (0, 0.2cm);
		}
	\end{tikzpicture}
	};
	\end{tikzpicture}
}

%New chapter command for book chapters
\newcommand{\GRPchapter}[1]{%

	\closeopenmulticols
	\begin{multicols}{2}[\chapter{#1}]\SetChapterTint}

%quick note formatting
\newenvironment{GRPquicknote}{}{}

\tcolorboxenvironment{GRPquicknote}{breakable,
halign=flush center,
leftrule=1.8mm,
rightrule=1.8mm,
toprule=0.5mm,
bottomrule=0.4mm,
enhanced jigsaw,
colback=\ChapterBackgroundTint!10!white,
opacityback=0.4,
colframe=\ChapterBackgroundTint!75!black,
overlay first and middle={%
\draw[draw=\ChapterBackgroundTint!75!black,line width=0.5mm,line cap=round,decorate,decoration={zigzag,segment length=2mm,amplitude=0.5mm}] ($(frame.south west)+(0.25mm,0)$) -- ($(frame.south east)-(0.25mm,0)$);
\fill[\ChapterBackgroundTint!75!black] (frame.south west) -- ($(frame.south west)+(0.9mm,-0.9mm)$) -- ($(frame.south west)+(1.8mm,0)$) -- cycle;
\fill[\ChapterBackgroundTint!75!black] (frame.south east) -- ($(frame.south east)+(-0.9mm,-0.9mm)$) -- ($(frame.south east)+(-1.8mm,0)$) -- cycle;
},
overlay middle and last={%
\draw[draw=\ChapterBackgroundTint!75!black,line width=0.5mm,line cap=round,decorate,decoration={zigzag,segment length=2mm,amplitude=0.5mm}] ($(frame.north west)+(1.55mm,0)$) -- ($(frame.north east)-(1.55mm,0)$);
\fill[\ChapterBackgroundTint!75!black] (frame.north west) -- ($(frame.north west)+(0.9mm,0.9mm)$) -- ($(frame.north west)+(1.8mm,0)$) -- cycle;
\fill[\ChapterBackgroundTint!75!black] (frame.north east) -- ($(frame.north east)+(-0.9mm,0.9mm)$) -- ($(frame.north east)+(-1.8mm,0)$) -- cycle;
},
}

\makeatletter
\newenvironment{GRPwidefig}
{%
	\SetChapterTint{}
	\par\medskip\noindent\begin{tcolorbox}[%
halign=flush center,
leftrule=1.8mm,
rightrule=1.8mm,
toprule=0.5mm,
bottomrule=0.4mm,
enhanced jigsaw,
colback=\ChapterBackgroundTint!10!white,
opacityback=0.4,
colframe=\ChapterBackgroundTint!75!black,]%
	\begin{minipage}{\textwidth}
	\captionsetup{font={bf,singlespacing,color=\ChapterBackgroundTint!10!white}}%
	\begingroup%
	\let\oldcaption\caption%
	\renewcommand{\caption}[2][]{%
		\def\@captype{figure}%
		\begin{tcolorbox}[colback=\ChapterBackgroundTint!75!black,colframe=\ChapterBackgroundTint!75!black,coltext=\ChapterBackgroundTint!10!white,size=title]%
		\caption@caption[##1]{##2}%
		\end{tcolorbox}%
	}%
}%
{%
	\let\caption\oldcaption%
	\endgroup%
	\end{minipage}%
	\end{tcolorbox}%
	\par\medskip%
}
\makeatother

\AddToHook{env/GRPwidefig/after}{\begin{multicols}{2}}
\AddToHook{env/GRPwidefig/before}{\end{multicols}}

\renewcommand{\thefigure}{\thesection-\alph{figure}}



